\origpage{45}{GRAPH MINORS. XIV}
\noindent is disconnected. But $r \subseteq \Lambda_2$, since $bd(\Lambda) \cap bd(\Lambda_2) = \emptyset$, and it is easy to see that this contradicts the maximality of $\Lambda$. \fix{We give the details. The proof of (6.3) uses its hypothesis only as follows: for a region $r$ of $H$ with $r \cap \Lambda = \emptyset$, distinct $v_1, v_2 \in V(C)$, and a line $F_2$ with ends $v_1, v_2$ and $F_2 - \{v_1, v_2\} \subseteq r$, it finds a line $L \subseteq \bar{r} \cap \bar{\Lambda}$ with ends $v_1, v_2$, that is, a path of $C$ between $v_1$ and $v_2$ drawn in $\bar{r}$. Hence we may choose $r$ so that there are such $v_1, v_2, F_2$, for which neither of the two paths $P_1, P_2$ of $C$ between $v_1$ and $v_2$ satisfies $U(P_i) \subseteq \bar{r}$. Then $r \subseteq \Lambda_2$ as above, and so $F_2 \subseteq \Lambda_2$. For $i = 1, 2$, let $D_i \subseteq \bar{\Lambda}_2$ be the closed disc bounded by $F_2 \cup U(P_i)$. Now $F_2 \not\subseteq \bar{\Lambda}$, since $\bar{r} \cap \Lambda = \emptyset$ and $r \cap U(C) = \emptyset$. Suppose that $U(P_1) \not\subseteq D_2$ and $U(P_2) \not\subseteq D_1$. Since also $F_2 \not\subseteq \bar{\Lambda}$, the closed discs $\bar{\Lambda}$ and $D_2$ meet only in $U(P_2)$, and their union is a closed disc in $\bar{\Lambda}_2$ bounded by $F_2 \cup U(P_1)$. So it equals $D_1$, by uniqueness, which is impossible since it includes $U(P_2)$. Hence we may assume that $U(P_1) \subseteq D_2$. Then $U(C) \subseteq D_2$, and so $\bar{\Lambda} \subseteq D_2$, since $\bar{\Lambda}$ is the unique closed disc in $\bar{\Lambda}_2$ bounded by $U(C)$. Moreover $D_2 \subseteq \Lambda_2$, since $\text{bd}(D_2) = F_2 \cup U(P_2) \subseteq \Lambda_2$.

Since $U(P_1) \not\subseteq \bar{r}$, some edge $e_0$ of $P_1$ satisfies $e_0 \not\subseteq \bar{r}$, and then $e_0 \cap \bar{r} = \emptyset$, since an edge that meets $\bar{r}$ is incident with $r$. Let $Z = \Lambda \cup e_0$. Then $Z$ is non-empty and arc-wise connected, since $\Lambda \subseteq Z \subseteq \bar{\Lambda}$ and $\bar{\Lambda}$ is a closed disc. Also $Z \cap \text{bd}(D_2) = \emptyset$, since $\Lambda$ is disjoint from $r$ and from $U(C)$, and $e_0$ is disjoint from $r$, from $V(H)$, and from the edges of $P_2$. Hence $Z \subseteq D_2 - \text{bd}(D_2)$. Let $s$ be a region of $H$ with $\bar{s} \cap Z \ne \emptyset$. Then $s \ne r$, since $\bar{r} \cap Z = \emptyset$, and so $s$ is disjoint from $\text{bd}(D_2) \subseteq r \cup U(H)$. Since $s$ is connected and meets the open set $D_2 - \text{bd}(D_2)$, it follows that $s \subseteq D_2 - \text{bd}(D_2)$. By (5.2), applied with $D_2$ in place of $\Delta$, there is a circuit $C'$ of $H$ with $U(C') \subseteq D_2$, bounding an open disc $\Lambda' \subseteq D_2$ including $Z$. Then $U(C') \subseteq \Lambda_2$, $\Lambda' \subseteq \Lambda_2$, and $\Lambda_1 \subseteq \Lambda \subseteq \Lambda'$. But $e_0 \subseteq \Lambda' - \Lambda$, and so $\Lambda' \ne \Lambda$, contrary to the maximality of $\Lambda$.}\fixnote{fix:9.3-step-added}{The paper says only that it is easy to see that $r$ contradicts the maximality of $\Lambda$. The argument above is the one used in the formalization.}

Thus $H \cap (\Sigma - \Lambda)$ is rigid, as required. $\square$

\stmt
(9.4) \textit{Let $H$ be a 2-cell drawing in $\Sigma$, let $z \in A(H)$, and let $\mathscr{T}$ be a respectful tangle in $H$ of order $\theta$, with metric $d$. Let $1 \le \kappa \le \theta - 7$, and let $\Lambda$ be a $\kappa$-zone around $z$, including every $x \in A(H)$ with $d(z, x) \le 3$. Let $H$ be a subgraph of $G$, let $\mathscr{S}$ be a $\mathscr{T}$-local 3-segregation of $G$, and let $\alpha$ be an arrangement of $\mathscr{S}$ in $\Sigma$ compatible with $H$. Then there is a closed disc $\Delta \subseteq \Sigma$ such that $\Lambda \subseteq \Delta$, $bd(\Delta)$ is $\alpha$-normal, $H \cap \Sigma - \Delta$ has a circuit, and $d(z, x) \le \kappa + 3$ for all $x \in A(H)$ with $x \cap \Delta \ne \emptyset$.}

\textit{Proof.} By (9.2) there is a $(\kappa + 3)$-zone $\Lambda_1$ around $z$ such that $x \subseteq \Lambda_1$ for every $x \in A(H)$ with $d(z, x) \le \kappa$. Then $\bar{\Lambda} \subseteq \Lambda_1$, since $\Lambda$ is a $\kappa$-zone. Let $C_1 = H \cap bd(\Lambda_1)$, $C_2 = H \cap bd(\Lambda)$. Now $d(z, x) \le \kappa + 3 \le \theta - 4$ for all $x \in A(H)$ with $x \subseteq \bar{\Lambda}_1$; and $x \subseteq \Lambda$ for all $x \in A(H)$ with $d(z, x) \le 3$. By (7.7), there is an $\alpha$-normal O-arc $F \subseteq \bar{\Lambda}_1$, bounding a closed disc $\Delta \subseteq \bar{\Lambda}_1$ with $\bar{\Lambda} \subseteq \Delta$. Then $\Delta$ satisfies the theorem. $\square$

Let $H$ be a 2-cell drawing in $\Sigma$, let $z \in A(H)$, and let $\mathscr{T}$ be a respectful tangle in $H$ of order $\theta$, with metric $d$. Let $\beta \ge 4$ be an even integer. A triple $(\Lambda, C, X)$ is a \textit{battlefield of size $\beta$ around $z$} if

\hspace*{12pt}(i) $\Lambda$ is a $(3\beta + 3)$-zone around $z$, and $H \cap (\Sigma - \Lambda)$ is a rigid drawing in $\Sigma$

\hspace*{12pt}(ii) $X \subseteq V(H) \cap bd(\Lambda)$, and $\theta \ge 12\beta + 15$, and $X$ is free with respect to the tangle of order $\theta - 4(3\beta + 3) - 2$ obtained by clearing $\Lambda$,

\hspace*{12pt}(iii) $C$ is a circuit of $H$ with $U(C) \subseteq \Lambda$, and there are $\beta$ disjoint paths of $H$ from $U(C)$ to $X$, each with no vertex in $H \cap bd(\Lambda)$ except its last,

\hspace*{12pt}(iv) the regional distance between $U(C)$ and $bd(\Lambda)$ is at least $1 + \frac{1}{2}\beta$,

\hspace*{12pt}(v) the open disc in $\Lambda$ bounded by $U(C)$ includes every $x \in A(\delm{K}\fix{H})$ with $d(z, x) < \beta$.\fixnote{fix:9-battlefield}{A misprint: no radial drawing $K$ appears in the definition, and $d$ is the metric on $A(H)$. The proofs of (9.5) and (9.6) check and use this clause for $x \in A(H)$.}

The main result of this section is the following.

\stmt
(9.5) \textit{Let $H$ be a 2-cell, rigid drawing in $\Sigma$, let $z \in A(H)$, and let $\mathscr{T}$ be a respectful tangle in $H$ of order $\theta$. Let $\beta \ge 4$ be an even integer with $\theta \ge 16\beta + 17$. Then there is a battlefield around $z$ of size $\beta$.}

\textit{Proof.} Let $d$ be the metric of $\mathscr{T}$. By two applications of (9.2), with $\kappa = \beta$ and $\kappa = 3\beta + 1$, we deduce that there is a $(\beta + 2)$-zone $\Lambda_2$ around $z$ such that $x \subseteq \Lambda_2$ for every $x \in A(H)$ with $d(z, x) < \beta$, and a $(3\beta + 3)$-zone $\Lambda_1$ around $z$ such that $x \subseteq \Lambda_1$ for every $x \in A(H)$ with $d(z, x) \le 3\beta$. Then $\bar{\Lambda}_2 \subseteq \Lambda_1$ because $\beta + 2 \le 3\beta$, and $z \subseteq \Lambda_2$ because $d(z, z) = 0 < \beta$.
